% ID: SEL-D0036-OA-E04
% ESTADO: cotejada. Segunda pasada visual y borrador compilado revisados.
% Fuente: Selectividad/matematicas_aplicadas/Ua1MATEMÁTICAS APLICADAS examen.pdf
% SHA-256: 77e0c295529ccd4ed502b5859d4011cb4b0e955a1e4795e570a9da20ae48293c
% Página: 1; opción A; ejercicio original 4.

\textbf{(2.5 puntos)} La talla media de los alumnos de una Universidad sigue una distribución Normal de media 170 cm y desviación típica 6 cm. Estudios recientes hacen sospechar que dicha talla media ha aumentado. Para confirmar, o no, esa sospecha se ha tomado una muestra de 64 estudiantes de esa Universidad, cuya talla media ha resultado ser de 172 cm. Con un nivel de significación del 1\%, plantee un contraste de hipótesis $(H_0:\mu\leq170)$, determine la región crítica de ese contraste y razone si se puede concluir que la talla media poblacional ha aumentado.
